\section{Introduction}
\label{sec:introduction}

Reproducibility in machine-learning research depends not only on the
availability of source code, but also on experimental choices encoded in that
code and its execution environment. Random seeds, data partitioning,
cross-validation protocols, dependency versions, and related configuration
choices can affect measured results even when the software continues to execute
successfully \cite{reimers2017score,henderson2018matters,bouthillier2021variance}.
Accordingly, reproducibility initiatives increasingly emphasize the
publication of code, environments, experimental details, and executable
workflows \cite{gundersen2018reproducibility,pineau2021improving,heil2021standards}.

An open question is whether the validation mechanisms already present in
research repositories are sensitive to changes in these
reproducibility-relevant choices. A repository may contain tests, continuous
integration, validation scripts, documented experiments, or runnable examples.
Such workflows provide evidence that software can execute or that particular
properties hold, but they are not necessarily designed to detect changes in
experimental configuration.

This distinction is difficult to study from ordinary repository outcomes
alone. When a validation workflow succeeds, it is generally unknown whether it
would also succeed after a scientifically relevant experimental choice had
changed. Mutation testing provides a way to make this counterfactual explicit:
introduce a controlled change and observe whether the existing validation
workflow distinguishes the mutant from the original program
\cite{jia2011mutation,just2014mutants}.

Mutation testing has already been applied extensively to conventional software,
scientific software, and machine-learning systems
\cite{hook2009trustworthiness,ma2018deepmutation,humbatova2021deepcrime}.
Dependency-update testing has likewise used artificial faults to assess whether
test suites can detect problems introduced through changed dependencies
\cite{hejderup2022dependencies}. The present work therefore does not claim that
mutation testing itself is new to ML or research software. Instead, we use
mutation testing for a specific measurement target: the sensitivity of
\emph{existing validation workflows in real ML research repositories} to
controlled changes in reproducibility-relevant experimental choices.

We introduce MLReproMutate, research software that detects and applies
domain-specific mutations to reproducibility-relevant choices in Python ML
projects. In the present study, we consider four mutation classes:
\texttt{random-seed}, which changes a literal random seed;
\texttt{dependency-pin}, which relaxes an exact dependency constraint;
\texttt{data-split}, which disables explicit stratification in a
train/test split; and \texttt{cv-fold-count}, which changes an explicit
cross-validation fold count. Each mutation is evaluated against a validation
workflow selected from the repository rather than against a newly introduced
study-specific test oracle.

We use MLReproMutate in an empirical study of 39 frozen
repository--operator cases drawn from real ML research repositories. Corpus
construction was static and outcome-blind: repository eligibility, source
revision, mutation candidate, and validation workflow were fixed before
mutation outcomes were observed. For the primary B02 corpus, validation-oracle
categories were also recorded prospectively. Cases that could not execute at
baseline were kept separate from mutation outcomes rather than being replaced
after observing execution results.

Primary execution produced mutation outcomes for 13 of the 39 cases. A bounded
post-primary restoration procedure increased the combined evaluable set to 24
cases while preserving the frozen revision, mutation candidate, and validation
workflow. One evaluated mutation was confirmed semantically equivalent. Among
the remaining 23 confirmed non-equivalent mutations, the selected validation
workflows detected 2 and did not detect 21. Thus, \emph{in this sample}, the
observed confirmed mutation-detection proportion was 2/23 (8.7\%).

We additionally examine detection descriptively by validation-workflow type and,
for the prospectively classified B02 subset, by oracle type. The two detected
B02 mutations both occurred in the
\texttt{dependency-pin} $\times$ \texttt{completion-only} cell. Sparse cells,
unequal operator composition, and differential evaluability prevent attributing
this pattern to oracle strength. We therefore treat these comparisons as
descriptive rather than as evidence of causal or statistical differences among
workflow categories.

The study makes four contributions:

\begin{itemize}
    \item We present a reproducibility-oriented mutation model and its
    implementation in MLReproMutate, targeting experimental choices surrounding
    ML research workflows rather than learned-model structure alone.

    \item We define an outcome-blind empirical protocol that freezes repository
    revisions, mutation candidates, validation workflows, and prospective oracle
    metadata before mutation outcomes are observed.

    \item We provide empirical evidence from 39 frozen
    repository--operator cases, including explicit accounting for baseline
    executability, bounded restoration, and semantic equivalence.

    \item We distinguish mutation survival from repository-level
    reproducibility: a surviving confirmed non-equivalent mutation establishes
    only that the selected validation workflow did not detect that particular
    controlled change.
\end{itemize}

The remainder of the paper reviews related work
(Section~\ref{sec:related-work}), describes MLReproMutate and its mutation model
(Section~\ref{sec:mlrepromutate}), states the research questions and empirical
protocol, reports the results, and discusses their implications and threats to
validity.
